\chapter{Results on the grinding circuit}
\label{ch:mill-results}
% Purpose (THESIS_GUIDE.md, Ch. 6): the answer to the research questions.
% Numbers (mill repo, lib/llm-supervisor/docs/): explore_exploit.md,
% explore_runs/fair_v1_explore_runs.md, explore_runs/reward_terms_v1.md,
% tdmpc/fair_v1*.md, pilot_300s/final_report.md, normalised_score.md,
% noise/noise_v1.md, feed_noise.md, noise_training.md,
% vm_week_2026-11-04_13.md. Sections 6.1-6.7 are exploratory; 6.8 is the
% confirmatory result of the registered protocol.
Sections~\ref{sec:mill-runs}--\ref{sec:mill-champions} are exploratory: their
champions were selected on the development and validation batteries they are
compared on, the method changed between runs, and there are at most three runs
per setting, with no significance test. They show what the method does and
chose the registered settings. The confirmatory comparison on the sealed test
battery is Section~\ref{sec:mill-sealed}. All results are at 300~s decisions
unless stated. Scores are means over a battery (lower is better), rewards the
TD-MPC's reward averaged over the battery (higher is better, at most 1), and
the number of terminal episodes is given in parentheses.

\section{Training runs: what the search found and what it cost}
\label{sec:mill-runs}
Table~\ref{tab:mill-runs} lists the training runs at 300~s, all with Claude
Opus 5.5 at effort low, three candidates per generation and 12 generations.
A run costs US\$3.2--3.9 and about 2.5~hours on one machine.

\begin{table}[tb]
  \centering
  \caption{Exploratory training runs at 300~s (12 generations, Opus 5.5,
    effort low). Development score of the final champion.}
  \label{tab:mill-runs}
  \small
  \begin{tabular}{@{}llrr@{}}
    \toprule
    Run & Prompt version & Cost (US\$) & Development \\
    \midrule
    \texttt{pilot\_300s}     & 5 (one champion) & 3.85 & 67.0 \\
    \texttt{protocol\_run1}  & 5 & -- & 70.0 \\
    \texttt{protocol\_run2}  & 5 & -- & 76.2 \\
    \texttt{explore\_run1}   & 6 (explore, exploit) & 3.40 & 69.5 \\
    \texttt{explore\_run2}   & 6 & 3.62 & 67.8 \\
    \texttt{explore\_run3}   & 6 & 3.23 & 67.9 \\
    \bottomrule
  \end{tabular}
\end{table}
\TODO{Fill the costs of protocol\_run1--2; add the eight registered protocol
runs when done. Mention the earlier 60~s runs (v4, mill
\texttt{docs/repeated\_runs\_v4.md}) in one sentence.}

\paragraph{How a run progresses.} The pilot's development score fell from
2\,754.7 (the fixed recipe) to 175.3 after the first generation and then in
eight smaller steps to 67.0. Of its 36
candidates, 9 were promoted, 15 rolled back by the guard, 11 not selected and 1
rejected on validation. After the first generation nearly every candidate kept
94--100\% of the champion's lines (median 98\%): a run polishes the design it
finds first.
\TODO{Figure: champion development score per
generation for every 300~s run (data in each run's \texttt{final\_report.md}
and \texttt{summary.csv}; plotting script to write).}

\paragraph{What the LLM found.} Every good champion uses the same levers: the
mill speed below 13~rpm for energy (lowering it from 13 to 12~rpm cuts mill
power by about 0.5~MW), the sump density setpoint to correct P80 and the sump
level together, and an opposite cyclone-pressure move to keep the level steady
while P80 is corrected. The lessons the LLM kept are statements about the
plant, for example that ``lighter ore both coarsens P80 and lowers the sump
level, so lowering sump density fixes both at once''.

\paragraph{Spread between runs.} With prompt version~5, runs with the same
settings ended far apart (development 67.0--76.2, reward 0.76--0.92), set by
how early they lowered the mill speed. Version~6 explores three lineages
(energy, level and quality priorities) for four generations before exploiting
the best; its three runs spread by 1.7 points on development against 6.2--9.2
for version~5, and every version~6 champion has a higher reward than both
version~5 protocol runs. The quality lineage was chosen in all three.
Continuing the version~6 runs to 16 generations gained 0.8--1.9\% for two of
them and nothing for the third, and an explore phase at high effort cost twice
as much without a gain. These results chose the registered settings: version~6,
12 generations, effort low.

\section{Main comparison at 300 s}
\label{sec:mill-main}
Table~\ref{tab:mill-main} compares the champions with the TD-MPC and the
baselines. The four best champions score lower than the TD-MPC on every
battery. On the TD-MPC's own reward, \texttt{pilot\_300s},
\texttt{explore\_run1} and \texttt{explore\_run2} are higher everywhere;
\texttt{explore\_run3} is slightly lower on development (0.876 against 0.887)
and validation (0.854 against 0.856), and the version~5 run
\texttt{protocol\_run1} is lower on every battery. The largest
differences are beyond the range and on validation, where the TD-MPC's score
rises to 111.6 against 66.8--76.3 for the champions. No champion has a
terminal episode on any of the three batteries; the TD-MPC has three beyond
the range.

\begin{table}[tb]
  \centering
  \caption{Main comparison at 300~s: score / reward / normalised score per
    battery (normalised score where computed). Exploratory: the champions were
    selected on these batteries.}
  \label{tab:mill-main}
  \footnotesize\setlength{\tabcolsep}{2pt}
  \begin{tabular}{@{}llll@{}}
    \toprule
    Controller & Development & Validation & Beyond \\
    \midrule
    \texttt{pilot\_300s}    & 67.0 / 0.922 / 0.904 & 66.8 / 0.920 / 0.905 & 194.1 / 0.725 / 0.810 \\
    \texttt{explore\_run1}  & 69.5 / 0.899         & 68.2 / 0.894         & 153.0 / 0.760 \\
    \texttt{explore\_run2}  & 67.8 / 0.931         & 67.0 / 0.924         & 151.0 / 0.770 \\
    \texttt{explore\_run3}  & 67.9 / 0.876         & 76.3 / 0.854         & 232.4 / 0.701 \\
    \texttt{protocol\_run1} & 70.0 / 0.761         & 78.0 / 0.764         & 238.2 / 0.605 \\
    TD-MPC                  & 72.5 / 0.887 / 0.881 & 111.6 / 0.856 / 0.854 & 662.7 / 0.656 / 0.704 (3) \\
    Decentralised PI        & 474.4 / 0.513 / 0.693 (1) & 189.9 / 0.583 / 0.755 & 1851.2 / 0.281 / 0.483 (7) \\
    Fixed recipe            & 2754.7 / 0.188 / 0.384 (18) & 2712.2 / 0.187 / 0.409 (16) & 4310.4 / 0.118 / 0.231 (21) \\
    \bottomrule
  \end{tabular}
\end{table}
\TODO{Normalised scores of the explore and protocol runs: run the
normalised report on the VM's stored episodes (no API calls).}
% python scripts/normalised_report.py runs/benchmark/fair_v1_explore_runs.csv --intervals 300

\subsection{By objective}
\label{sec:mill-perspectives}
Table~\ref{tab:mill-terms} splits the development score into its objectives.
Energy is 63.5--65.0 points of every good controller's score (6.35--6.50
kWh/t) and the differences between them are small; the TD-MPC uses slightly
more. The control penalty, what is left without energy, is 3.3--4.7 for the
champions against 7.5 for the TD-MPC, mostly from the sump-level band (4.64
weighted points against 0.6--1.0). On validation the TD-MPC's level-band term
rises to 35.6 and its P80 term to 3.7, against at most 8.4 and 0.5 for the
champions. Beyond the range, the level band and safety dominate every
controller's control penalty.

\begin{table}[tb]
  \centering
  \caption{Development score split by objective (weighted points, so the
    columns add up to the score), and the specific energy.}
  \label{tab:mill-terms}
  \small\setlength{\tabcolsep}{3pt}
  \begin{tabular}{@{}lrrrrrrrr@{}}
    \toprule
    Controller & Score & P80 & Level & CL band & Safety & Energy & CL target & kWh/t \\
    \midrule
    \texttt{pilot\_300s}   & 67.0 & 0.50 & 0.64 & 0.93 & 0 & 63.70 & 1.23 & 6.37 \\
    \texttt{explore\_run1} & 69.5 & 1.26 & 1.04 & 0.93 & 0 & 64.85 & 1.47 & 6.49 \\
    \texttt{explore\_run2} & 67.8 & 1.50 & 0.75 & 0.93 & 0 & 63.53 & 1.14 & 6.35 \\
    \texttt{explore\_run3} & 67.9 & 0.65 & 0.68 & 0.93 & 0 & 64.36 & 1.26 & 6.44 \\
    TD-MPC                 & 72.5 & 0.48 & 4.64 & 0.93 & 0 & 65.02 & 1.47 & 6.50 \\
    Decentralised PI       & 474.4 & 9.58 & 154.04 & 10.14 & 228.57 & 69.29 & 2.80 & 6.93 \\
    \bottomrule
  \end{tabular}
\end{table}
\TODO{Same split for validation and beyond (in the source docs), or in an
appendix. Add P80 IAE from the target and setpoint travel as reported,
unscored measures (mill \texttt{docs/tdmpc/fair\_v1\_metrics.md}).}

\subsection{Where the differences are}
\label{sec:mill-where}
The TD-MPC was trained on constant feed, and its gap is in the scenarios that
vary the feed over time. On validation it scores 71.2 on the steady and step
scenarios, about as the champions (68.1--70.4), but 142.0 on ramps and sines,
against 65.4--80.7. Beyond the range its ramp and sine score is 892.3.
\TODO{Per-scenario paired differences (figure: one dot per scenario,
champion minus TD-MPC, score and reward). One trajectory figure: a
development scenario, champion against TD-MPC (level, P80, setpoints).}

\section{Beyond the development range}
\label{sec:mill-beyond}
From development to beyond, the champions' scores rise by a factor of
2.2--3.4 and their rewards fall by 0.14--0.20; the TD-MPC's score rises by a
factor of 9.1 and its reward falls by 0.23, with three terminal episodes. What
fails is mostly the sump-level band (65.9--132.8 weighted points for the
champions beyond the range, 86.1 for the TD-MPC) and then P80; the TD-MPC also
pays 454.6 points of safety from its terminal episodes, the champions none.
\TODO{Per kind and pattern. The degradation against amplitude (1, 1.25, 1.5,
2) comes from the sealed battery (Section~\ref{sec:mill-sealed}).}

\subsection{Failure modes}
\label{sec:mill-failure-modes}
\TODO{Catalogue of failure modes (C2), each with one example trajectory, and
an empirical operating envelope for the LLM supervisor.}

\section{Score and reward disagree on the sump level}
\label{sec:mill-reward}
The score charges nothing while the sump level is inside 70--80\%. The
reward's level factor is a narrow bell around 75\%: 1\% off costs about 8\% of
its multiplicative part, 2\% off about 26\%. Champions with almost the same
score therefore differ in reward by how tightly they hold 75\%
(Table~\ref{tab:mill-reward-terms}): on development \texttt{explore\_run2}
holds the level 0.6\% from 75\% on average and loses 0.058 reward to the level,
\texttt{explore\_run3} 1.0\% and 0.095, and the version~5 run
\texttt{protocol\_run1} 1.3\% and 0.142. Beyond the range the level is the
largest reward loss for every champion. The trainer promotes on the score, so
this is not selected for. This explains every case in
Table~\ref{tab:mill-main} where a champion wins on score and loses on reward.

\begin{table}[tb]
  \centering
  \caption{The TD-MPC's reward split by term on the development battery:
    reward gained if that term alone were perfect, and the mean distance of
    the sump level from 75\%.}
  \label{tab:mill-reward-terms}
  \small
  \begin{tabular}{@{}lrrrrr@{}}
    \toprule
    Controller & Reward & P80 & Level & CL & Level off 75\% \\
    \midrule
    \texttt{pilot\_300s}    & 0.922 & 0.024 & 0.062 & 0.047 & 0.7\% \\
    \texttt{explore\_run1}  & 0.899 & 0.023 & 0.062 & 0.061 & 0.8\% \\
    \texttt{explore\_run2}  & 0.931 & 0.020 & 0.058 & 0.045 & 0.6\% \\
    \texttt{explore\_run3}  & 0.876 & 0.027 & 0.095 & 0.050 & 1.0\% \\
    \texttt{protocol\_run1} & 0.761 & 0.042 & 0.142 & 0.079 & 1.3\% \\
    \bottomrule
  \end{tabular}
\end{table}
\TODO{Add the TD-MPC's row (run \texttt{scripts/reward\_terms.py} for it) and
a figure of the reward split per controller.}

\section{Decision interval}
\label{sec:mill-interval}
Every controller does best at the interval it was developed for
(Table~\ref{tab:mill-interval}). Champions trained at 60~s beat the TD-MPC at
60~s and degrade at 300~s; the champion trained at 300~s beats it at 300~s and
900~s and stays usable at 60~s. At 900~s the TD-MPC falls to near the level of
the decentralised PI, with seven terminal episodes on development.

\begin{table}[tb]
  \centering
  \caption{Development battery at three decision intervals: score / reward
    (terminal episodes). \texttt{v4\_low\_run1} is the best of three champions
    trained at 60~s.}
  \label{tab:mill-interval}
  \small\setlength{\tabcolsep}{4pt}
  \begin{tabular}{@{}llll@{}}
    \toprule
    Controller & 60~s & 300~s & 900~s \\
    \midrule
    \texttt{pilot\_300s} (trained at 300~s) & 89.0 / 0.833 & 67.0 / 0.922 & 441.9 / 0.678 (1) \\
    \texttt{v4\_low\_run1} (trained at 60~s) & 65.3 / 0.980 & 82.9 / 0.885 & 580.2 / 0.651 (3) \\
    TD-MPC (trained at 300~s)              & 69.3 / 0.922 & 72.5 / 0.887 & 1119.6 / 0.585 (7) \\
    Decentralised PI                       & 184.4 / 0.605 (1) & 474.4 / 0.513 (1) & 1182.2 / 0.460 (5) \\
    \bottomrule
  \end{tabular}
\end{table}
\TODO{Validation and beyond columns, or a second table; the other two 60~s
champions (development 65.0 and 63.5 at 60~s, 125.1 and 331.7 at 300~s).}

\section{Robustness to noise}
\label{sec:mill-noise-results}
Table~\ref{tab:mill-noise-results} gives the development reward under each
noise category of Section~\ref{sec:mill-noise}, and
Table~\ref{tab:mill-noise-beyond} the terminal episodes beyond the range.
Under feed fluctuations every champion keeps a higher reward than the TD-MPC on
every battery; under measurement noise \texttt{explore\_run2} falls below it.

\begin{table}[tb]
  \centering
  \caption{Development reward by noise category (one draw per scenario).
    \texttt{noise\_run1--2} were trained under plant measurement noise (prompt
    version~7) and are tested on fresh draws.}
  \label{tab:mill-noise-results}
  \small
  \begin{tabular}{@{}lrrrr@{}}
    \toprule
    Controller & None & Measurement & Feed & Feed + measurement \\
    \midrule
    \texttt{explore\_run1} & 0.899 & 0.682 & 0.692 & 0.586 \\
    \texttt{explore\_run2} & 0.931 & 0.529 & 0.695 & 0.480 \\
    \texttt{explore\_run3} & 0.876 & 0.782 & 0.658 & 0.620 \\
    \texttt{noise\_run1}   & 0.801 & 0.757 & 0.629 & 0.614 \\
    \texttt{noise\_run2}   & 0.890 & 0.779 & 0.677 & 0.641 \\
    TD-MPC                 & 0.887 & 0.653 & 0.610 & 0.500 \\
    Decentralised PI       & 0.513 & 0.481 & 0.322 & 0.300 \\
    \bottomrule
  \end{tabular}
\end{table}

\begin{table}[tb]
  \centering
  \caption{Terminal episodes on the beyond battery (24 scenarios) by noise
    category.}
  \label{tab:mill-noise-beyond}
  \small
  \begin{tabular}{@{}lrrrr@{}}
    \toprule
    Controller & None & Measurement & Feed & Feed + measurement \\
    \midrule
    \texttt{explore\_run1} & 0 & 1 & 0 & 5 \\
    \texttt{explore\_run2} & 0 & 1 & 0 & 2 \\
    \texttt{explore\_run3} & 0 & 1 & 0 & 1 \\
    TD-MPC                 & 3 & 6 & 7 & 11 \\
    Decentralised PI       & 7 & 10 & 11 & 16 \\
    \bottomrule
  \end{tabular}
\end{table}

\subsection{Measurement noise on P80}
The TD-MPC does not read the measured P80, so P80 noise leaves it unchanged.
The champion \texttt{pilot\_300s} barely moves: development 67.0 / 0.922
without noise, 68.2 / 0.914 with Gaussian noise of 5~$\mu$m and 71.1 / 0.892
with the analyser model.

\subsection{Measurement noise on every signal}
Plant-wide measurement noise reverses the noise-free ranking of the champions.
\texttt{explore\_run2}, the best without noise (reward 0.931), becomes the
worst champion (0.529), and \texttt{explore\_run3}, the weakest, the most
robust (0.782). The cause in \texttt{explore\_run2} is one line: it estimates
the sump-level trend from two samples 60~s apart. Replacing only that line with
a 5-minute least-squares slope lifts its development reward under noise from
0.529 to 0.706 and lowers it without noise from 0.931 to 0.834. Selecting on
noise-free scenarios does not select for robustness, as on the tank test bed
(Section~\ref{sec:dev-lessons}).

\subsection{Feed fluctuations}
Fluctuating feed hurts every controller, the TD-MPC most: its development
reward falls from 0.887 to 0.610, and its terminal episodes beyond the range
from 3 to 7, against none for the champions. Under feed fluctuations alone the
champions trained without noise do best. Adding measurement noise on top
separates them again, as without feed fluctuations.

\subsection{Training under noise}
Two runs trained under plant measurement noise (prompt version~7, outside the
protocol) became insensitive to it, but not clearly better than the most robust
noise-free champion, and training on measurement noise bought nothing against
feed fluctuations.

\subsection{Measuring under noise}
One noise draw per scenario is not enough for a close comparison: one
scenario's score doubled between draws (66.6 to 140.3). A champion trained
under noise must be tested on fresh draws.
\TODO{Fairness caveat in one sentence: the TD-MPC was trained without noise,
and Weir's feed generator could randomise its training feed. An amplitude
sweep within one category would make RQ2's noise axis a curve
(Section~\ref{sec:mill-noise}).}

\section{The generated supervisors}
\label{sec:mill-champions}
A champion is a single deterministic Python function. \texttt{pilot\_300s} has
125 lines; the version~6 champions have 62--85. The pilot's structure is a
level-trend estimator that drives the cyclone pressure setpoint, an
ore-density feedforward and a P80 trim on the sump density setpoint, a rule for
the number of cyclones, and the mill speed held near 11.8~rpm for energy. The
full listing is in Appendix~\ref{app:mill-champion}.
\TODO{Rule count per champion (the goal document's complexity measure: count
branches and thresholds; cheap to script). Readability, honestly: the mill
champions use terse names and few comments. The \texttt{explore\_run2}
diagnosis above is the best evidence that the code can be inspected and fixed.}

\section{Confirmatory comparison on the sealed test battery}
\label{sec:mill-sealed}
% Filled in after the eight protocol runs and the one sealed evaluation
% (Section sec:protocol). All tests at alpha = 0.05.
\TODO{After the eight protocol runs. If they are not run, this section becomes
a list of what the exploratory results can and cannot claim (champions
selected on development and validation; three runs per setting; no
significance test; validation in the promotion rule; TD-MPC trained on
constant feed for another reward).}

\subsection{$H_{0,1}$: within the development range}
\TODO{Mean paired difference LLM minus TD-MPC on \texttt{test\_dev} with 95\%
intervals, score and reward (co-primary); equivalence (TOST, 90\% interval
within $\pm 3.5$ score and $\pm 0.02$ reward); Cohen's $d_z$ with its interval,
descriptive.}

\subsection{$H_{0,2}$: degradation beyond the range}
\TODO{Slope of mean $\log(1 + J_s)$ against amplitude level 1, 1.25, 1.5, 2,
LLM minus TD-MPC, with 95\% interval; figure of score and reward against
amplitude for each controller; the unseen sine periods.}

\subsection{C3: spread between independent runs}
\TODO{Per-policy means (median, IQR, range), between-policy SD with a
chi-square interval, runs without a promotion.}

\subsection{Secondary analyses}
\TODO{Per objective (P80 outside range, P80 IAE, level band, CL band, safety,
travel, energy, CL target) with Holm-corrected $p$-values; per kind and
pattern; steady/step against ramp/sine; the 60~s, 900~s and noise variants.}

\section{Summary}
\label{sec:mill-summary}
\TODO{Rewrite after the sealed test. Exploratory so far: an LLM-written
deterministic supervisor of 60--130 lines matches or beats the deployed TD-MPC
on this simulator at its own decision interval, on both measures for most
runs, and fails more gracefully beyond the range and under noise (fewer
terminal episodes). The result depends on the run (spread), the measure (score
against reward on the sump level) and the conditions (noise reverses the
ranking among champions).}
